\documentclass[12pt]{article}
\usepackage[utf8]{inputenc}
\usepackage[T1]{fontenc}
\usepackage{hyperref}
\title{Pao Language Resources for K-12 Curriculum Development}
\author{Compiled by Claude Sonnet 5.5 (claude-sonnet-5-5) for Andr\'as Kornai}
\date{}
\begin{document}
\maketitle

\section{Introduction}
The code ppa (Pao) is not a living language entry: Glottolog lists it as retired from ISO 639-3, merged into Bhageli [bfy] (change request 2015-012, adopted 2016-01-15), on the grounds that the Pao people speak a variety of Bundelkhandi and the separate entry was based on a misunderstanding.

\section{Primary Educational Resources}
None; there is no distinct Pao language to which they could belong.

\section{Linguistic and Reference Materials}
\begin{itemize}
\item Glottolog record: \url{https://glottolog.org/resource/languoid/id/paoo1245}
\item ISO 639 ``p'' code listing: \url{https://en.wikipedia.org/wiki/ISO_639:p}
\end{itemize}

\section{Religious and Community Resources}
None found.

\section{Digital and Audio Resources}
None found.

\section{Curriculum Development Context}
No curriculum can be built for this code. Work should instead target the language the community actually speaks, which according to Glottolog is a variety of Bundelkhandi, and the retired-code status should be recorded.
\end{document}